\chapter{AI Harness Tools: Mastering the Agentic Developer Toolkit}

\begin{quote}
\textit{``In 2023, the best AI coding tool was a chat window. By mid-2025, developers have access to over a dozen autonomous coding agents --- some run in IDEs, some in terminals, some in fully sandboxed cloud environments. They all share a common architecture: a ReAct loop, a context engine, a diff applicator, and a permission model. The tools differ in how they implement each of these components, and understanding those differences is the key to choosing the right tool for your workflow.''}
\end{quote}

This chapter provides a deep-dive architectural analysis of every major AI coding tool in the modern agentic ecosystem. For each tool, we investigate its internal architecture, context indexing mechanisms, diff application strategies, and permission designs.

\section{The Architectural Taxonomy of AI Coding Tools}

Before examining individual implementations, we classify the ecosystem into three architectural paradigms:

\begin{enumerate}
    \item \textbf{IDE-Integrated Clients (e.g., Cursor, Windsurf, Copilot):} Deeply integrated with the editor UI. They leverage editor state (active cursor position, open files, diagnostics, git diffs) to assemble context, and apply modifications directly to the editor buffer.
    \item \textbf{CLI-Native Terminal Agents (e.g., Claude Code, Aider):} Execute directly in the developer's terminal. They are git-native, highly procedural, and rely on standard command-line tools to read, edit, test, and commit code.
    \item \textbf{Sandboxed Cloud SWE Agents (e.g., Devin, OpenHands):} Run entirely in remote sandboxes (Docker/MicroVMs). They have full browser control, terminal terminals, and are designed for asynchronous task delegation.
\end{enumerate}

\section{Detailed Tool Architecture \& Internals}

\subsection{1. Cursor: Deep IDE Integration \& Context Synthesis}

Cursor is a fork of VS Code. Instead of running as an extension, running as a separate IDE allows Cursor to modify the core editor loop to enable autocomplete and inline diff rendering.

\begin{figure}[ht]
\centering
\begin{tikzpicture}[
    scale=0.74, transform shape,
    node distance=0.6cm and 1cm,
    box/.style={draw, rectangle, rounded corners, minimum height=0.7cm, minimum width=2.8cm, align=center, fill=blue!5, font=\small},
    arrow/.style={thick,->,>=stealth}
]
\node[box] (edit) {Editor State \\ (Buffer, Selection, Linter)};
\node[box, right=of edit] (context) {Context Engine \\ (AST Parser + BM25)};
\node[box, below=of context] (server) {Cursor Backend \\ Orchestrator};
\node[box, below=of server] (llm) {Speculative LLM \\ (Diff Plan)};
\node[box, left=of llm] (myers) {Myers Diff Engine};
\node[box, above=of myers] (editor) {Editor Buffer \\ Inline Diff};

\draw[arrow] (edit) -- (context);
\draw[arrow] (context) -- (server);
\draw[arrow] (server) -- (llm);
\draw[arrow] (llm) -- (myers);
\draw[arrow] (myers) -- (editor);
\draw[arrow] (editor) -- (edit);
\end{tikzpicture}
\caption{Cursor Backend Processing Architecture.}
\end{figure}

\subsubsection{Speculative Autocomplete Engine}
Cursor's inline autocomplete does not use standard Claude or GPT models directly due to latency requirements ($<100$ms). Instead, it runs custom speculative models. The local model suggests token sequences based on the current editor context, and a larger cloud-hosted model validates these proposals. This allows multi-line predictions (e.g., writing the entire function body) to stream instantly.

\subsubsection{Context Indexing (AST + BM25)}
Cursor builds a local index of the codebase using tree-sitter. It parses symbols (classes, methods, variables) into an Abstract Syntax Tree (AST) (Chapter 6). When the user invokes a query, the orchestrator combines:
\begin{itemize}
    \item \textbf{Semantic Similarity:} Cosine similarity on embeddings of AST nodes.
    \item \textbf{Fuzzy Match (BM25):} Traditional lexical search to match exact variable/class names.
    \item \textbf{Scope Analysis:} Evaluating import hierarchies to pull in dependent definitions.
\end{itemize}

\subsubsection{Myers Diff Engine}
When applying edits (Cmd-K or Agent Mode), Cursor avoids rewriting the entire file. Instead, the model outputs raw edits, and Cursor uses an optimized variant of the \textit{Myers Diff Algorithm} (Chapter 11) to compute line-level insertions and deletions, rendering them as green/red inline review blocks.

\subsection{2. Claude Code: CLI-Native Thinking \& Executing Agent}

Claude Code represents a shift from IDE GUI to procedural CLI. It is a pure ReAct agent designed to run command-line tools locally.

\subsubsection{Extended Thinking \& Chain of Thought}
Claude Code is optimized to leverage Anthropic's Claude models with extended thinking. Before executing any tool call, the model outputs thousands of internal thinking tokens to simulate code outcomes, write pseudocode, and verify syntax.

\subsubsection{The Permission Enforcement Loop}
Because Claude Code runs locally on the user's host machine (not in a sandbox), security is critical. It operates a strict client-side verification loop:
\begin{itemize}
    \item \textbf{Read Commands:} Automatic permission. Reading files or checking git status occurs silently.
    \item \textbf{Write Commands:} File writes are inspected. If the write modifies files outside the project root, it blocks.
    \item \textbf{System Execution (Bash):} Every command (e.g., \texttt{npm test}, \texttt{git commit}) is shown to the user with a confirmation prompt. The user can whitelist safe commands (like \texttt{pytest}) for the current session.
\end{itemize}

\begin{lstlisting}[language=Bash]
# Typical Claude Code run loop execution
Thinking: The user wants to run tests. I need to call the bash tool.
Proposed Command: npm run test:unit
[Claude Code CLI Prompts User]: Run 'npm run test:unit'? (Y/n)
\end{lstlisting}

\subsection{3. Aider: Git-Native Pair Programming \& Repository Mapping}

Aider is an open-source terminal pair programming tool. Its primary innovations are in context size reduction and git integration.

\subsubsection{Repository Mapping via Ctags}
To query large codebases without exceeding token limits, Aider generates a \textbf{Repository Map}. It uses \texttt{universal-ctags} to extract symbols (class definitions, function signatures) from all files. This map is represented as a tree structure:
\begin{lstlisting}[language=Python]
# Aider repository map output representation
src/auth/helper.py:
    class AuthHelper:
        def validate_token(self, token: str) -> bool:
        def refresh_session(self) -> str:
\end{lstlisting}
This map is injected into the system prompt, giving the LLM a structural index of the repository. When the LLM needs to edit, it queries this map to identify which files to load into full context.

\subsubsection{Architect vs. Editor Pattern}
Aider uses a split-model workflow:
\begin{enumerate}
    \item \textbf{Architect Mode:} A reasoning-heavy model (e.g., OpenAI o1/o3-mini) receives the user query and generates a detailed implementation plan.
    \item \textbf{Editor Mode:} A fast, diff-proficient model (e.g., Claude 3.5 Sonnet) receives the plan and applies the actual code modifications to the files.
\end{enumerate}

\subsection{4. Devin \& OpenHands: Sandboxed Multi-Agent Systems}

Devin and OpenHands are autonomous SWE agents that run inside isolated cloud sandboxes.

\subsubsection{Sandbox Isolation Architecture}
To prevent malicious code or accidental infinite loops from compromising host systems, the agent executes within a secure Docker container or MicroVM running custom overlay filesystems (Chapter 8). The agent has access to a sandboxed shell, a sandboxed web browser, and a local directory structure.

\subsubsection{Event-Driven Multi-Agent Collaboration}
OpenHands uses an event-driven loop centered around a shared \textbf{Event Stream}. The orchestrator dispatches tasks to specialized sub-agents:
\begin{itemize}
    \item \textbf{Planner Agent:} Sets high-level objectives.
    \item \textbf{Coder Agent:} Writes and edits files.
    \item \textbf{Executor Agent:} Runs tests in the sandboxed shell.
    \item \textbf{Browser Agent:} Interacts with web pages using Playwright/Puppeteer.
\end{itemize}
Every action (run command, modify file, view web page) and observation (terminal output, linter warnings, web page screenshots) is published as an event. Agents subscribe to this stream and react asynchronously.

\begin{figure}[ht]
\centering
\begin{tikzpicture}[
    scale=0.76, transform shape,
    node distance=0.7cm and 1cm,
    box/.style={draw, rectangle, rounded corners, minimum height=0.7cm, minimum width=2.5cm, align=center, fill=blue!5, font=\small},
    agentbox/.style={draw, rectangle, rounded corners, minimum height=0.7cm, minimum width=2.2cm, align=center, fill=yellow!15, font=\small},
    arrow/.style={thick,->,>=stealth}
]
% Central Stream
\node[box, fill=purple!10, minimum width=8cm] (stream) {Shared Event Stream \\ (Actions, Observations, Comments)};

% Agents
\node[agentbox, above=of stream, xshift=-3cm] (planner) {Planner Agent};
\node[agentbox, above=of stream, xshift=-1cm] (coder) {Coder Agent};
\node[agentbox, above=of stream, xshift=1.5cm] (executor) {Executor Agent};
\node[agentbox, above=of stream, xshift=3.8cm] (browser) {Browser Agent};

% Sandbox Tools
\node[box, below=of stream, xshift=-2cm] (shell) {Sandboxed Shell};
\node[box, below=of stream, xshift=2cm] (playwright) {Playwright Browser};

% Connections
\foreach \a in {planner, coder, executor, browser} {
    \draw[arrow] (\a) -- (stream);
    \draw[arrow] (stream) -- (\a);
}
\draw[arrow] (stream) -- (shell);
\draw[arrow] (shell) -- (stream);
\draw[arrow] (stream) -- (playwright);
\draw[arrow] (playwright) -- (stream);

\end{tikzpicture}
\caption{Event-Driven Sandbox Agent Architecture (OpenHands / Devin model).}
\end{figure}

\section{Unified System Architecture Pattern}

While each tool is tailored for a specific workflow (terminal, IDE, or browser), they all share a universal core loop. 

\begin{figure}[ht]
\centering
\begin{tikzpicture}[
    scale=0.75, transform shape,
    node distance=0.8cm and 1.2cm,
    box/.style={draw, rectangle, rounded corners, minimum height=0.7cm, minimum width=3.5cm, align=center, fill=blue!8, font=\small},
    arrow/.style={thick,->,>=stealth}
]
\node (start) [box, fill=gray!15] {User Task / Query};
\node (context) [box, below=of start] {Context Gatherer \\ (Fuzzy search, AST symbols, Git state)};
\node (planner) [box, below=of context] {LLM Orchestrator / Planner \\ (Generates sub-tasks \& plans)};
\node (diff) [box, below=of planner] {Diff applicator \\ (Line-level updates / AST validation)};
\node (sandbox) [box, below=of diff] {Tool Execution Sandbox \\ (Terminal runner, linter check)};
\node (eval) [box, right=3.5cm of planner] {Observation Evaluator \\ (Parse errors, diff verify)};

\draw[arrow] (start) -- (context);
\draw[arrow] (context) -- (planner);
\draw[arrow] (planner) -- (diff);
\draw[arrow] (diff) -- (sandbox);
\draw[arrow] (sandbox.east) -- node[right, font=\footnotesize, xshift=0.2cm] {Result} ++(2.5,0) |- (eval.south);
\draw[arrow] (eval.north) -- ++(0,1) -| (planner.east);
\end{tikzpicture}
\caption{Unified Agent Loop Architecture.}
\end{figure}

\section{Ecosystem Comparison Matrix}

The table below summarizes the key trade-offs across the agentic developer tools:

\begin{table}[h!]
\centering
\footnotesize
\begin{tabular}{|l|p{2.2cm}|c|c|p{2.8cm}|p{2.5cm}|}
\hline
\textbf{Tool} & \textbf{Paradigm} & \textbf{OSS} & \textbf{MCP} & \textbf{Primary Models} & \textbf{Host Env} \\
\hline
\textbf{Cursor} & IDE-Integrated & \ding{55} & \ding{51} & Claude Sonnet, GPT-4o & Local Host \\
\textbf{Claude Code} & CLI Agent & \ding{55} & \ding{51} & Claude 3.7 Sonnet & Local Host \\
\textbf{Copilot} & IDE Extension & \ding{55} & \ding{51} & GPT-4o, Claude Sonnet & Local / Cloud \\
\textbf{Windsurf} & IDE-Integrated & \ding{55} & \ding{51} & Cascade Custom, Claude & Local Host \\
\textbf{Aider} & CLI pair-prog & \ding{51} & \ding{55} & Any (BYO Key) & Local Host \\
\textbf{Roo Code} & VS Code Ext & \ding{51} & \ding{51} & Any (BYO Key) & Local Host \\
\textbf{Cline} & VS Code Ext & \ding{51} & \ding{51} & Any (BYO Key) & Local Host \\
\textbf{OpenHands} & SWE Agent & \ding{51} & \ding{55} & Any (BYO Key) & Sandboxed Docker \\
\textbf{Devin} & Cloud Agent & \ding{55} & \ding{55} & Proprietary reasoning & Cloud MicroVM \\
\hline
\end{tabular}
\caption{Comparative analysis of developer agent tools.}
\end{table}

\begin{center}
\noindent\rule{0.6\textwidth}{0.4pt} \\
\vspace{0.3em}
\textit{[Code implementation is available in the companion coding handbook: \texttt{coding-handbook/ch20\_ai\_harness\_tools/}]} \\
\noindent\rule{0.6\textwidth}{0.4pt}
\end{center}
